\documentclass[11pt,a4paper,sans]{moderncv}

% Estilo do ModernCV
\moderncvstyle{classic} 
\moderncvcolor{blue}

\usepackage[utf8]{inputenc}
\usepackage[scale=0.8]{geometry}
\usepackage{setspace} \setstretch{1.3}

% Dados Pessoais
\name{Fabio Eduardo de O.}{Medeiros}
\title{Doutor em Física | Data Scientist | Reliability Engineer}

\address{Fortaleza/CE}{}
\phone[mobile]{+55~(85)~98827-5798}
\email{feomedeirosdev@gmail.com}
\social[linkedin]{fabio-medeiros-708567262}
\social[github]{feomedeirosdev}
\extrainfo{38 anos | Solteiro}

\begin{document}
\makecvtitle

\section{Resumo Profissional}
Cientista de Dados e Doutor em Física da Matéria Condensada com mais de 10 anos de experiência em pesquisa aplicada, modelagem computacional e análise estatística. Especialista no desenvolvimento de frameworks para gestão de ativos (ULRA) e análise de custo de ciclo de vida (LCCA), unindo rigor científico e inteligência artificial para transformar comportamentos físicos de degradação em indicadores financeiros estratégicos. Experiência sólida em Python, estatística aplicada e aprendizado de máquina.

\section{Experiência Profissional e Acadêmica}

\cventry{2025--Atual}{Estrategista de Dados (Projeto LCCA)}{Desenvolvimento Independente / Framework ULRA}{}{}{Desenvolvimento de um Modelo de Otimização Estocástica para Análise de Custo de Ciclo de Vida (LCCA).
\begin{itemize}
    \item Implementação de algoritmos para classificação de perfis de degradação (\textit{Stable Heroes, Fast Decay, Bad Actors}).
    \item Análise de confiabilidade através de TTFF e modelagem de ciclos de vida via processos estocásticos.
    \item Desenvolvimento de matrizes de decisão para otimização de CAPEX/OPEX.
    \item \textbf{Tecnologias:} Python (Pandas, Scipy, Seaborn), Reliability Engineering, Mathematical Optimization.
\end{itemize}}

\cventry{Jun/2023--Jul/2024}{Cientista de Dados | Inovação no Judiciário}{Bolsa de Inovação Tecnológica (BIT)}{}{}{Atuação no projeto ``Ciência de Dados e IA em Apoio às Atividades Judiciárias''.
\begin{itemize}
    \item Modelagem de redes e aplicação de técnicas de grafos para análise de acórdãos e jurisprudências.
    \item Desenvolvimento de buscadores semânticos baseados em grafos para identificação de padrões legais.
    \item \textbf{Tecnologias:} Python, NetworkX, Pandas, Bibliotecas de Grafos.
\end{itemize}}

\cventry{Mar/2019--Fev/2025}{Pesquisador de Doutorado}{UFC}{Fortaleza}{}{Pesquisa em perovskitas fotovoltaicas com foco em propriedades ópticas e estruturais.
\begin{itemize}
    \item Modelagem computacional e análise estatística aplicada à caracterização experimental.
    \item Produção científica e orientação de novos pesquisadores.
    \item \textbf{Ferramentas:} Python, Origin, Fityk, ShellScript.
\end{itemize}}

\cventry{Mar/2018--Fev/2019}{Pesquisador de Mestrado}{UFC}{Fortaleza}{}{Caracterização de materiais ferróicos e multiferróicos, com análise de termoestabilidade.
\begin{itemize}
    \item Coleta de dados, interpretação de resultados e suporte em aulas de laboratório.
    \item \textbf{Ferramentas:} Python, Origin, Fityk, ShellScript.
\end{itemize}}

\cventry{Mar/2011--Mar/2014}{Bolsista de Iniciação Científica}{Grupo de Espectroscopia (UFC)}{}{}{Pesquisa experimental focada na caracterização de compostos multiferroicos e co-cristalização de fármacos.
\begin{itemize}
    \item Trabalho interdisciplinar envolvendo física, química e ciência de materiais.
\end{itemize}}

\cventry{2015--2016}{Professor de Física Experimental}{Sistema Farias Brito}{}{}{Preparação de estudantes para competições olímpicas (OBF, IPhO, IYPT). Representante da comitiva brasileira na OIbF 2015 e IYPT 2016.}

\section{Formação Acadêmica}
\cvitem{2019--2025}{\textbf{Doutorado em Física} -- Universidade Federal do Ceará (UFC)}
\cvitem{2018--2019}{\textbf{Mestrado em Física} -- Universidade Federal do Ceará (UFC)}
\cvitem{2009--2013}{\textbf{Bacharelado em Física} -- Universidade Federal do Ceará (UFC)}

\section{Competências Técnicas}
\cvitem{Linguagens}{Python, SQL, HTML5, CSS3, JavaScript, PHP, ShellScript, \TeX}
\cvitem{Data Science}{ML, LCCA, Modelagem Estocástica, Engenharia de Confiabilidade, Grafos}
\cvitem{Bibliotecas}{Pandas, Scikit-learn, Scipy, Reliability, NetworkX, Seaborn}
\cvitem{Ferramentas}{Git/GitHub (GitFlow), VS Code, Linux (Mint/Ubuntu)}

\vfill
\centerline{\footnotesize \copyright~2026 Fabio Medeiros - Físico \& Estrategista de Dados.}

\end{document}